%% kapa: exam=1 round=32 year=2021 period=1 form=A subject=economics subject_ko=경제학원론 seq=28 number=68
%% source: Q-Net 감정평가사 기출문제 (cst00309) artlSeq=5208397 file=extracted/1008/5208397/A형/2021년도 제32회 감정평가사 1차 1교시 A형.pdf
%% conversion: kordoc 4.24.51 -> tools/gen1.py (auto); review=proofread-sonnet; flags=none
\begin{question}{68}
실업률과 인플레이션율의 관계는 $u=u_n-2(\pi-\pi_e)$이고 자연실업률이 3\%이다. 보기를 고려하여 중앙은행이 0\%의 인플레이션율을 유지하는 준칙적 통화정책을 사용했을 때의 (ㄱ)\uline{실업률}과, 최적 인플레이션율로 통제했을 때의 (ㄴ)\uline{실업률은}? (단, $u$, $u_n$, $\pi$, $\pi_e$는 각각 실업률, 자연실업률, 인플레이션율, 기대 인플레이션율이다.)
\begin{bogi}
○ 중앙은행은 물가를 완전하게 통제할 수 있다. \par
○ 민간은 합리적인 기대를 하며 중앙은행이 결정한 인플레이션율로 기대 인플레이션율을 결정한다. \par
○ 주어진 기대 인플레이션에서 중앙은행의 최적 인플레이션율은 1\%이다. \par
\end{bogi}
\begin{choices}
  \item ㄱ: 0\%, ㄴ: 0\%
  \item ㄱ: 1\%, ㄴ: 0\%
  \item ㄱ: 1\%, ㄴ: 1\%
  \item ㄱ: 2\%, ㄴ: 1\%
  \item ㄱ: 3\%, ㄴ: 3\%
\end{choices}
\end{question}
\answer{5}
